
\section{Appendix A. Proof of Theorem 1 (Eigengap-Identification-HAC chain)}

We collect notation. Let $D\subset\mathbb{R}^d$ be a bounded open set with
Lipschitz boundary and Lebesgue volume $|D|<\infty$; let $\mu$ be Lebesgue
measure restricted to $D$. The locations $s_1,\dots,s_n\in D$ are deterministic
in the conditional argument below (we condition on $\{s_i\}$ throughout, as
in Gilbert--Ogburn--Datta 2024, Section 3.2). The Matern kernel of smoothness
$\nu>0$ and range $\rho>0$ is
\[
K_{\nu,\rho}(s,s') \;=\; \frac{1}{\Gamma(\nu)2^{\nu-1}}
   \!\Bigl(\tfrac{\sqrt{2\nu}\,\|s-s'\|}{\rho}\Bigr)^{\!\nu}
   \mathcal{K}_\nu\!\Bigl(\tfrac{\sqrt{2\nu}\,\|s-s'\|}{\rho}\Bigr),
\]
with $\mathcal{K}_\nu$ the modified Bessel function of the second kind. Its
spectral density on $\mathbb{R}^d$ satisfies
\begin{equation}\label{eq:matern-specdens}
\widehat{K}_{\nu,\rho}(\omega) \;=\;
   C_{\nu,\rho,d}\bigl(2\nu/\rho^2 + \|\omega\|^2\bigr)^{-(\nu+d/2)},
\qquad
C_{\nu,\rho,d}=\frac{\Gamma(\nu+d/2)(2\nu)^{\nu}}{\pi^{d/2}\Gamma(\nu)\,\rho^{2\nu}},
\end{equation}
(Stein 1999, eq.~(2.10); Rasmussen and Williams 2006, eq.~(4.15)). In
particular $\widehat{K}_{\nu,\rho}(\omega)\asymp\|\omega\|^{-(2\nu+d)}$ as
$\|\omega\|\to\infty$. We write $a_k\asymp b_k$ for ``$0<\liminf a_k/b_k\le
\limsup a_k/b_k<\infty$''.

We let $K_n=[K_{\nu,\rho}(s_i,s_j)]_{i,j=1}^n$ and $S_n=K_n(K_n+\lambda I_n)^{-1}$,
with $\lambda>0$ the smoothing parameter. The continuum integral operator
$\mathcal{T}_K:L^2(D,\mu)\to L^2(D,\mu)$ is
$(\mathcal{T}_K f)(s)=\int_D K_{\nu,\rho}(s,t)f(t)\,d\mu(t)$.

\medskip
\begin{theorem}\label{thm:main}
Under assumptions (A1)--(A4) the following five conclusions hold simultaneously,
with constants made explicit in their respective lemmas.
\begin{enumerate}
\item[\emph{(i)}] $\lambda_k(S_n)=\omega_k/(\omega_k+\lambda)$ where
   $\omega_k=\lambda_k(K_n)$; the bulk continuum eigenvalues satisfy
   $\omega_k\asymp k^{-(2\nu+d)/d}$, and \emph{provided} the BBP cutoff lies
   in the Widom polynomial tail, i.e.\ $\lambda\ll\rho^{2\nu+d}$ (equivalently
   $r=\lfloor\lambda^{-d/(2\nu+d)}\rfloor\gg\rho^{-d}$; see Lemma~\ref{lem:eigengap}),
   the eigengap of $S_n$ at the cutoff satisfies
   $\delta_r\asymp\lambda^{d/(2\nu+d)}$, and the integrated tail-bias of the
   smoother satisfies $\sum_{k>r}\lambda_k(S_n)\asymp\lambda^{2\nu/(2\nu+d)}$.
   When the cutoff falls in the pre-asymptotic plateau ($r\lesssim\rho^{-d}$),
   $\delta_r\asymp\lambda^{1/a_{\mathrm{eff}}(r)}$ with $a_{\mathrm{eff}}(r)$
   the local log-log slope of $\omega_k$ at $k\!\approx\!r$; this slower rate
   propagates into the constants of (ii)--(iii) but does not change the rates.
\item[\emph{(ii)}] With $\widehat\Sigma_T=n^{-1}T^\top T$ and $\widehat v_1$ its
   leading unit eigenvector, with probability at least $1-2\exp(-cn)$,
\[
\|\sin\Theta(\widehat v_1,v_1)\|_{\mathrm{op}}
   \le \frac{2^{3/2}\,\|\widehat\Sigma_T-\Sigma_T\|_{\mathrm{op}}}{\delta_{\mathrm{eff}}}
   \le \frac{C_1\sqrt{(p+\log(2/\eta))/n}}{\delta_{\mathrm{eff}}},
\]
where $\delta_{\mathrm{eff}}=\lambda_1(\Sigma_T)-\lambda_2(\Sigma_T)=
\mathrm{Var}(S_n U)$ and $C_1=c_0(\sigma_E^2+\mathrm{Var}(S_n U))$ with $c_0$
absolute (from Vershynin 2018, Thm.~4.7.1).
\item[\emph{(iii)}] The PC1--OLS estimator $\widehat\theta_{\mathrm{PC1}}=
   (\widehat v_1^\top T^\top T\widehat v_1)^{-1}\widehat v_1^\top T^\top Y$ has
\[
\bigl|\mathbb{E}[\widehat\theta_{\mathrm{PC1}}\mid T,S_nU]-\langle\beta,v_1\rangle\bigr|
   \;\le\; C_2\,\|\sin\Theta(\widehat v_1,v_1)\|_{\mathrm{op}}\,
   \bigl(\|\beta_\perp\|+\|g(s)\|_{2,n}\bigr)\,
   \sqrt{\mathrm{Var}(S_nU)/\sigma_E^2},
\]
with $C_2=2(1+o_p(1))$, where $\|g\|_{2,n}^2:=n^{-1}\sum_i g(s_i)^2$.
\item[\emph{(iv)}] Let $\theta_0=\langle\beta_\perp,v_\perp\rangle$ for any fixed
   unit $v_\perp\in\mathrm{range}(P_\perp)$. The cross-fitted DML partialling-out
   estimator $\widehat\theta_\perp$ satisfies
\[
\sqrt n\,(\widehat\theta_\perp-\theta_0)\;\xrightarrow{d}\;\mathcal N(0,\sigma_\perp^2),
\qquad
\sigma_\perp^2 = \frac{\mathbb{E}\bigl[\psi(W;\theta_0,\eta_0)^2\bigr]}
                {\bigl(\mathbb{E}[(P_\perp T v_\perp-\mathbb E[P_\perp T v_\perp\mid s])^2]\bigr)^{2}},
\]
provided the product nuisance rate $\|\widehat f-f_0\|_{L^2}\cdot\|\widehat g-g_0\|_{L^2}=o_p(n^{-1/2})$
holds (Chernozhukov et al.\ 2018, Thm.~4.1, eq.~(4.5)).
\item[\emph{(v)}] The fold-centered jackknife $\widehat V_{\rm JK}$ of
   Salerno--Wu--McCormick 2026, Thm.~1, satisfies $n\widehat V_{\rm JK}\xrightarrow{p}\sigma_\perp^2$.
\end{enumerate}
\end{theorem}



\begin{proof}[Proof of Theorem~\ref{thm:main}]
\textbf{(i)} Apply Lemma~\ref{lem:smoother-eig} to $K_n$ to get
$\lambda_k(S_n)=\omega_k/(\omega_k+\lambda)$. By Lemma~\ref{lem:matern-mercer}
the empirical $\omega_k$ inherit the continuum rate $\omega_k\asymp k^{-(2\nu+d)/d}$
up to constants that depend on $n$ via the discrete-continuum convergence
(absorbed). Lemma~\ref{lem:eigengap} gives the eigengap rate at the cutoff
$r=\lfloor\lambda^{-d/(2\nu+d)}\rfloor$ as $\delta_r\asymp\lambda^{d/(2\nu+d)}$,
and the bias-direction (tail) rate $\sum_{k>r}\lambda_k(S_n)\asymp\lambda^{2\nu/(2\nu+d)}$
that drives (iii).

\textbf{(ii)} Apply Lemma~\ref{lem:vershynin} to the rows of $T/\sqrt{n}$, which
are sub-Gaussian by (A4) and the sub-Gaussianity verification at the end of
Lemma~\ref{lem:vershynin}. With probability $\ge 1-2e^{-cn}$,
$\|\widehat\Sigma_T-\Sigma_T\|_{\mathrm{op}}\le C_1\sqrt{(p+\log(2/\eta))/n}$.
Plug into Lemma~\ref{lem:yws} with $r=1$ and $\delta_{\mathrm{eff}}=\mathrm{Var}(S_nU)$
to get the displayed sin-$\Theta$ bound. The bound is informative whenever
$\lambda\ll\omega_1$ (super-critical BBP).

\textbf{(iii)} Lemma~\ref{lem:pc1-bias} delivers the bias decomposition; the
dominant term is $\langle\beta_\perp,\widehat v_1-v_1\rangle$, bounded by
$\|\beta_\perp\|\sqrt 2\|\sin\Theta(\widehat v_1,v_1)\|_{\mathrm{op}}$. Adding
the $g$-bias and SNR-amplification $\sqrt{\mathrm{Var}(S_nU)/\sigma_E^2}$
gives Theorem~\ref{thm:main}(iii).

\textbf{(iv)} Apply Lemma~\ref{lem:god} to $P_\perp T$: the residual satisfies
GOD-2024 (A3.7), so $\beta_\perp$ is identified. Apply Lemma~\ref{lem:dml} with
$D=(P_\perp T)v_\perp$ to get the $\sqrt n$-CLT. The product rate
$\|\widehat m-m_0\|_{L^2}\cdot\|\widehat g-g_0\|_{L^2}=o_p(n^{-1/2})$ holds for
kernel-ridge nuisances at Matern-$\nu$ smoothness: the minimax rate is
$n^{-(2\nu+d)/(4\nu+2d)}$ (Tsybakov 2009 Cor.~2.3); for $\nu=2.5,d=2$ that is
$n^{-1/2}$, so the product is $o(n^{-1/2})$ comfortably.

\textbf{(v)} Apply Lemma~\ref{lem:swm}: $n\widehat V_{\mathrm{JK}}\to_p\sigma_\perp^2$,
giving asymptotically valid Wald intervals.
\end{proof}
